% Generated from locale/tr/content/computability/recursive-functions/composition.tex; do not hand-edit.


\olfileid[tr]{cmp}{rec}{com}
\olsection{Bileşke}

$f$ ve $g$ doğal sayıların iki tek yerli fonksiyonuysa bunların
bileşkesini alabiliriz: $h(x)=f(g(x))$. Yeni $h(x)$ fonksiyonu $f$ ve
$g$ fonksiyonlarından \emph{bileşke} yoluyla tanımlanır. Bunu birden
çok girdili fonksiyonlara genellemek istiyoruz.

Bunu yapmanın bir yolu şudur: $f$, $k$ yerli; $g_0,\dots,g_{k-1}$ ise
hepsi $n$ yerli olan $k$ fonksiyon olsun. Yeni bir $n$ yerli $h$
fonksiyonunu şöyle tanımlayabiliriz:
\[
h(x_0, \dots, x_{n-1}) =
f(g_0(x_0, \dots, x_{n-1}), \dots, g_{k-1}(x_0, \dots, x_{n-1})).
\]
$f$ ve bütün $g_i$ fonksiyonları hesaplanabilirse $h$ de
hesaplanabilirdir: $h(x_0,\dots,x_{n-1})$ hesaplamak için önce her
$i=0,\dots,k-1$ için $y_i=g_i(x_0,\dots,x_{n-1})$ değerlerini
hesaplayın. Ardından bu değerleri $f$'ye girdi verip
$h(x_0,\dots,x_{n-1})=f(y_0,\dots,y_{k-1})$ hesaplayın.

Bu, var olan fonksiyonları kullanarak yeni bir fonksiyon hesaplarken
olanları gereğinden fazla kısıtlayan bir niteleme gibi görünebilir.
Öncelikle bazen bir fonksiyonun bütün girdilerini kullanmayız; örneğin
$\Add$ fonksiyonunun ilkel özyinelemeli tanımında kullanılmak üzere
$g(x,y,z)=\Succ(z)$ tanımladık. Şu fonksiyonları kullanmamıza izin
verildiğini varsayalım:
\[
\Proj{n}{i}(x_0, \dots, x_{n-1}) = x_i.
\]
$\Proj{k}{i}$ fonksiyonlarına \emph{izdüşüm} fonksiyonları denir;
$\Proj{n}{i}$, $n$ yerli bir fonksiyondur. O zaman $g$ şöyle
tanımlanabilir:
\[
g(x, y, z) = \Succ(\Proj{3}{2}(x, y, z)).
\]
Burada $f$ rolünü tek yerli $\Succ$ oynar; dolayısıyla $k=1$.
$g_0$ rolünü oynayan tek bir üç yerli $\Proj{3}{2}$ fonksiyonumuz
vardır. Sonuç, üçüncü girdinin ardılını döndüren üç yerli bir
fonksiyondur.

İzdüşüm fonksiyonları, girdileri yeniden sıralayarak veya özdeşleyerek
yeni fonksiyonlar tanımlamamıza da olanak verir. Örneğin
$h(x)=\Add(x,x)$ fonksiyonu
\[
h(x_0)=\Add(\Proj{1}{0}(x_0),\Proj{1}{0}(x_0))
\]
ile tanımlanabilir. Burada $k=2$, $n=1$; $f(y_0,y_1)$ rolünü $\Add$,
$g_0(x_0)$ ve $g_1(x_0)$ rollerinin ikisini de tek yerli izdüşüm
fonksiyonu (başka bir deyişle özdeşlik fonksiyonu)
$\Proj{1}{0}(x_0)$ oynar.

$f(y_0,y_1)$ elimizde bulunan bir fonksiyonsa
$h(x_0,x_1)=f(x_1,x_0)$ fonksiyonunu
\[
h(x_0, x_1) = f(\Proj{2}{1}(x_0, x_1),\Proj{2}{0}(x_0, x_1))
\]
ile tanımlayabiliriz. Burada $k=2$, $n=2$; $g_0$ ve $g_1$ rollerini
sırasıyla $\Proj{2}{1}$ ve $\Proj{2}{0}$ oynar.

$g_0,\dots,g_{k-1}$ fonksiyonlarının tümünün aynı $n$ sayıda girdiye
sahip olması da kaygı uyandırabilir. (Bir fonksiyonun \emph{yer sayısı}
girdi sayısıdır; $n$ yerli bir fonksiyonun yer sayısı $n$'dir.) Ancak
izdüşüm fonksiyonlarını eklemek istenen esnekliği sağlar. Örneğin $f$
ve $g$ üç yerli, $h$ ise
\[
h(x,y)=f(x,g(x,x,y),y)
\]
ile tanımlanan iki yerli bir fonksiyon olsun. $h$ tanımı izdüşüm
fonksiyonlarıyla şöyle yeniden yazılabilir:
\[
h(x,y) = f(\Proj{2}{0}(x,y), g(\Proj{2}{0}(x,y), \Proj{2}{0}(x,y),
\Proj{2}{1}(x,y)), \Proj{2}{1}(x,y)).
\]
O zaman $h$, $f$ fonksiyonunun $\Proj{2}{0}$, $l$ ve $\Proj{2}{1}$
ile bileşkesidir; burada
\[
l(x,y)=g(\Proj{2}{0}(x,y),\Proj{2}{0}(x,y),\Proj{2}{1}(x,y)),
\]
yani $l$, $g$ fonksiyonunun $\Proj{2}{0}$, $\Proj{2}{0}$ ve
$\Proj{2}{1}$ ile bileşkesidir.

